\documentclass[10pt]{article}
\usepackage[margin=0.8in]{geometry}
\usepackage[T1]{fontenc}
\sloppy
\begin{document}
\section*{Independent branch norm transport review}
Source fidelity supported; kernel unverified. Five prospective declarations were inspected against the actual imported maps and prior independent branch derivation. No compiler or fresh axiom audit was run.
\paragraph{Binding.} Producer revision \texttt{332575b279b0}; exact candidate bytes matched the Git object. Preread plan preceded source inspection. The reviewer did not author this candidate; announced claims and earlier derivation were known.
\paragraph{Actual maps.} The evaluator $E$ uses the existing projective coefficient homomorphism and $X$ coordinate. Its equality with affine polynomial evaluation followed by localization and the inverse branch-stalk equivalence is checked on constants and $X$. The scalar tower supplies coefficient compatibility.
\paragraph{Curve and norm.} The quotient equation yields $E(f)=Y^2$. Mapping the polynomial identity gives
\[N^2-E(u_1)E(u_2)E(U)=B^2Y^2,\]
and hence $(N-BY)(N+BY)=E(u_1)E(u_2)E(U)$. The factor $B^2$ comes from mapping $C(b^2)$ through both ring homomorphisms. The target norm equation is derived, not assumed.
\paragraph{Orders.} At a simple branch point, nonzero $b$ and support vanishing imply $n(a)=0$. The predecessor gives each conjugate order one. Applying additive valuation gives
\[v_P(E(u_1))+v_P(E(u_2))+v_P(E(U))=2.\]
The source specialization derives the polynomial norm from the coefficient expressions and divisor congruences, and derives the simple-root condition from separability.
\paragraph{Adversarial checks.} Zero $b$ and singular branches are excluded for the order claim. A zero support polynomial would contradict the finite product order, so it cannot satisfy all premises. Characteristic two does not invalidate these local identities, but distinct-conjugate or separable-cover interpretations must not be added there. The evaluator is tied to the actual stalk; no desired norm or order is a wrapper premise.
\paragraph{Precise remaining work.} The sum two is a stalk valuation statement. The independent hand derivation gives base-polynomial multiplicity sum one; that transport is not yet implemented. Exact compilation with complete pinned imports and independent five-name axiom readback remain required. The source file's print commands are not executed evidence. Global divisor and infinity obligations are separate.
\end{document}
